\documentclass[12pt, a4paper]{article}

% =============================================
% PAKETEAK
% =============================================
\usepackage[utf8]{inputenc}
\usepackage[T1]{fontenc}
\usepackage[basque]{babel}
\usepackage{geometry}
\usepackage{graphicx}
\usepackage{booktabs}
\usepackage{svg}
\usepackage{array}
\usepackage{float}
\usepackage{hyperref}
\usepackage{xcolor}
\usepackage{amsmath}
\usepackage{amssymb}
\usepackage{fancyhdr}
\usepackage{parskip}
\usepackage{caption}
\usepackage{multirow}
\usepackage{longtable}
\usepackage{pgfplots}
\pgfplotsset{compat=1.18}

\geometry{top=2.5cm, bottom=2.5cm, left=3cm, right=3cm}

\hypersetup{colorlinks=true, linkcolor=black, citecolor=black, urlcolor=blue}

\pagestyle{fancy}
\fancyhf{}
\rhead{Spam Sailkatzailea -- MLP}
\lhead{Erabakiak Hartzeko Euskarri Sistemak}
\cfoot{\thepage}

% =============================================
% DOKUMENTUA
% =============================================
\begin{document}

% =============================================
% AZALA
% =============================================
\begin{titlepage}
    \centering
    \vspace*{1.5cm}
    {\Large \textbf{BILBOKO INGENIARITZA ESKOLA}}\\[0.3cm]
    {\normalsize Kudeaketaren eta Informazio Sistemen Informatikaren Ingeniaritzarako Gradua}\\[2cm]
    \rule{\textwidth}{1.5pt}\\[0.5cm]
    {\Huge \textbf{Spam Sailkatzailea}}\\[0.4cm]
    {\Large Multilayer Perceptron }\\[0.3cm]
    \rule{\textwidth}{1.5pt}\\[2cm]

    \begin{minipage}{0.45\textwidth}
        \textbf{Egileak:}\\
        Gorka Piedra\\
        Andoni Ortiz de Zarate\\
        Ibai Olaziregi\\
    \end{minipage}
    \begin{minipage}{0.45\textwidth}
        \raggedleft
        \textbf{Arloa:} Erabakiak Hartzeko Euskarri Sistemak\\[0.2cm]
        \textbf{Maila:} 3. maila\\
        \textbf{Lauhilabetea:} 2. lauhilabetea\\
    \end{minipage}
    \vfill
    {\large \today}
\end{titlepage}

% =============================================
% AURKIBIDEA
% =============================================
\tableofcontents
\newpage

% =============================================
% 1. TALDE-LANAREN ARDUREN BANAKETA
% =============================================
\section{Talde-lanaren arduren banaketa}

\begin{table}[H]
    \centering
    \caption{Taldekide bakoitzaren ardura eta inplementatutako zatiak}
    \begin{tabular}{p{2.5cm}p{4cm}p{6.5cm}}
        \toprule
        \textbf{Kidea} & \textbf{Ardura nagusia} & \textbf{Inplementatutako zatiak} \\
        \midrule
        Gorka Piedra &
        Dokumentazio teknikoa, kalitatearen ebaluazioa eta datuen prestaketa. &
        \textbullet{} JavaDoc dokumentazioaren sorrera proiektuko klase anitzetan. \newline
        \textbullet{} Datuak partitzeko logikaren inplementazioa (\textit{DataSplit}). \newline
        \textbullet{} Ereduaren kalitate-probak egitea. \newline
        \textbullet{} Erroreen kudeaketa eta \texttt{doku.tex} dokumentazioaren fintzea. \newline
        \textbullet{} Fluxu-diagramaren diseinua. \\
        \midrule
        Andoni Ortiz de Zarate &
        Esperimentazio azpiegitura, entrenamendua eta ereduaren iraunkortasuna. &
        \textbullet{} MLP (\textit{Multilayer Perceptron}) garatzea eta doitzea \textit{grid search} bidez. \newline
        \textbullet{} Esperimentuen erregistroa kudeatzea \texttt{Evaluate.java} eta \texttt{ExperimentLogger.java} fitxategietan. \newline
        \textbullet{} \texttt{GetModel.java} logika eredu entrenatua lortzeko eta gordetzeko. \newline
        \textbullet{} Datuen hasierako karga eta ereduaren testeatzea. \\
        \midrule
        Ibai Olaziregi &
        Testuaren aurreprozesamendua, iragarpen-interfazea eta kodearen garbiketa. &
        \textbullet{} \texttt{EmailVectorizerAndSelector.java} garatzea bektorizaziorako eta atributu-hautapenerako. \newline
        \textbullet{} \texttt{Iragarri.java} inplementatzea mezu berrien iragarpenerako. \newline
        \textbullet{} Karaktere berezien tratamendua eta normalizazioa \texttt{TextNormalizer.java} fitxategian. \newline
        \textbullet{} Klase-diagramaren diseinua eta proiektuaren azken antolaketa (errefaktorizazioa eta garbiketa). \\
        \bottomrule
    \end{tabular}
    \label{tab:ardura}
\end{table}

Aipatzekoa da taldekide guztiok parte hartu dugula, neurri handiagoan edo txikiagoan, proiektuaren atal guztietan (hala nola dokumentazioan, kodearen garbiketetan edota testaketan). Hala ere, taulan agertzen den banaketak atal bakoitzean pisu handiena izan duen edo lan karga handiena bere gain hartu duen kidea islatzen du, ardura nagusiak argiago identifikatzeko asmoz.

\newpage
% =============================================
% 2. DATU-SORTAREN ANALISIA
% =============================================
\section{Datu-sortaren analisia}

\subsection{Datu gordinak}

Erabilitako datu-multzoa Enron-Spam corpusa da. Bi klasetan banatutako
mezu elektronikoz osatuta dago: \textbf{spam} (nahi ez dena) eta \textbf{ham} (legitimo).

Spam mezuak sustapen komertzialen, iruzurren eta premium zerbitzuen ingurukoak dira,
hizkera limurtzailea, maiuskulak, presazko estiloa eta estekak erabiliz
(adib.: \textit{``Free prize! Call now to claim your \pounds 1000''}).

Ham mezuak eguneroko komunikazio pertsonalak dira, hizkera informala,
laburdurak eta emotikonoak erabiliz (adib.: \textit{``Hey, what are you doing?''}).

\begin{table}[H]
    \centering
    \caption{Jatorrizko datu-multzoaren deskribapen kuantitatiboa (datu gordinak)}
    \begin{tabular}{lrrr}
        \toprule
        \textbf{Ezaugarria} & \textbf{Train} & \textbf{Dev} & \textbf{Test} \\
        \midrule
        Atributu nominalak  & 1     & 1     & 1     \\
        Atributu numerikoak & 0     & 0     & 0     \\
        String atributuak   & 1     & 1     & 1     \\
        Atributu guztira    & 2     & 2     & 2     \\
        \midrule
        Instantzia (spam $\oplus$)  & 1.200 & 150   & 150   \\
        Instantzia (ham $\ominus$)  & 2.937 & 368   & 367   \\
        Instantzia guztira          & 4.137 & 518   & 517   \\
        \bottomrule
    \end{tabular}
    \label{tab:datuak_gordinak}
\end{table}

Taula~\ref{tab:datuak_gordinak}-n ikusten den bezala, klaseen artean desoreka
nabarmena dago: spam instantziak ham instantzien \%29 baino ez dira.
Horrek eragina du ereduaren errendimenduan, bereziki spam klaseko recall-ean.

\subsection{Partiketa estratifikatua}

\texttt{DataSplit.java} fitxategiak \texttt{StratifiedRemoveFolds} metodoa erabiliz
datu-multzoa hiru azpimultzotan banatzen du, klaseen proportzioa mantenduz
partiketa guztietan. Partiketa \textbf{\%80/\%10/\%10} egin da:

\begin{itemize}
    \item \textbf{Train (\%80):} 4.137 instantzia -- eredua entrenatzeko
    \item \textbf{Dev (\%10):} 518 instantzia -- parametroak optimizatzeko
    \item \textbf{Test (\%10):} 517 instantzia -- azken iragarpenak egiteko
\end{itemize}

\newpage
% =============================================
% 3. AURREPROZESAMENDUA
% =============================================
\section{Aurreprozesamendua}

\subsection{Aplikatutako prozesuak}

\texttt{EmailVectorizerAndSelector.java} fitxategiak \texttt{StringToWordVector}
filtroaren bidez testu gordinak bektore numeriko bihurtzen ditu.
Prozesu honetan honako eragiketak egin dira:

\begin{itemize}
    \item \textbf{Minuskulizazioa (\texttt{LowerCaseTokens=true}):}
    Maiuskulak eta minuskulak berdintzen dira. \textit{``FREE''} eta \textit{``free''}
    hitz bera bezala tratatzen dira.
    \item \textbf{Tokenizazioa (\texttt{WordTokenizer}):}
    Testua hitzetan banatzen da puntuazio-markak banatzailetzat erabiliz. \\
    Muga-karaktereak: \texttt{" \textbackslash r\textbackslash n\textbackslash t.,;:'``()?!-+/\textbackslash<>@\#\$\%\^{}\&*\_=\textasciitilde`|[]\{\}"}
    \item \textbf{Stemming (\texttt{LovinsStemmer}):}
    Hitzak beren erroetara murrizten dira. Adib.: \textit{``running''}, \textit{``runs''}
    eta \textit{``run''} hitz bera bilakatzen dira.
    \item \textbf{Stopwords (\texttt{Rainbow}):}
    Hitz ez-informatiboak ezabatzen dira (\textit{``the''}, \textit{``is''}, \textit{``at''}, ...).
\end{itemize}

\subsection{Parametroen eragina (kuantitatiboa)}

Stemmer-aren eta stopword-en eraginaren ebaluazioa train multzoan egin da.
\textbf{TODO: Osatu esperimentu honekin -- stemmer bai/ez konparaketa egin eta
taula bete.}

\begin{table}[H]
    \centering
    \caption{Stemmer-aren eraginaren konparaketa (F-Measure, spam)}
    \begin{tabular}{lrr}
        \toprule
        \textbf{Konfigurazioa} & \textbf{Atributu kopurua} & \textbf{F-spam} \\
        \midrule
        Stemmer gabe  & -- & -- \\
        LovinsStemmer & -- & -- \\
        \bottomrule
    \end{tabular}
    \label{tab:stemmer}
\end{table}

\newpage
% =============================================
% 4. BEKTORIZAZIOA
% =============================================
\section{Bektorizazioa}

\subsection{Bektorizazio-metodo ezberdinak}

Lau bektorizazio-metodo aztertu dira train multzoan, MLP sailkatzailearekin batera,
egokiena aukeratzeko. Ebaluazio-eskema: \textbf{Hold-Out} (train vs. dev).
Metrika: \textbf{F-Measure (spam)}.

\begin{table}[H]
    \centering
    \caption{Bektorizazio metodoen konparaketa (F-Measure)}
    \begin{tabular}{lrrr}
        \toprule
        \textbf{Metodoa} & \textbf{F-spam} & \textbf{F-ham} & \textbf{WAvg-F} \\
        \midrule
        TF-IDF   & -- & -- & -- \\
        TF       & -- & -- & -- \\
        IDF      & -- & -- & -- \\
        Bin BoW  & -- & -- & -- \\
        \bottomrule
    \end{tabular}
    \label{tab:bektorizazio}
\end{table}

\textbf{TODO: Osatu taula bektorizazio-metodo bakoitzarekin esperimentatu ondoren.}

\subsection{StringToWordVector parametroak}

\begin{table}[H]
    \centering
    \caption{Aukeratutako StringToWordVector parametro guztiak}
    \begin{tabular}{lll}
        \toprule
        \textbf{Parametroa} & \textbf{Balioa} & \textbf{Zergatia} \\
        \midrule
        TFTransform     & true   & Hitz maiztasuna kontuan hartzeko \\
        IDFTransform    & true   & Hitz arraroak gehiago ponderatzeko \\
        LowerCaseTokens & true   & Maiuskulak berdintze \\
        WordsToKeep     & 20.000 & Hiztegi zabala mantentzeko atributu-hautaketarako \\
        Tokenizadorea   & WordTokenizer & Puntuazio-markak banatzailetzat \\
        Stemmer         & LovinsStemmer & Hitzen erroetara murrizteko \\
        Stopwords       & Rainbow & Hitz ez-informatiboak kentzeko \\
        \bottomrule
    \end{tabular}
    \label{tab:stwv}
\end{table}

\subsection{Atributuen hautapena}

\texttt{AttributeSelection} filtro bat aplikatu da \texttt{InfoGainAttributeEval}
ebaluatzailearekin eta \texttt{Ranker} bilatzailearekin. Informazio-irabaziak
spam klasearekiko erlazionatuen dauden atributuak aukeratzen ditu.

\textbf{Hautaketaren eragina:}

\begin{table}[H]
    \centering
    \caption{Atributu-hautaketaren eragina}
    \begin{tabular}{lrr}
        \toprule
        \textbf{Konfigurazioa} & \textbf{Atributu kopurua} & \textbf{F-spam} \\
        \midrule
        Hautaketa gabe (WordsToKeep=20.000) & 20.000 & -- \\
        Hautaketarekin (Ranker=1.000)       & 1.000  & -- \\
        \bottomrule
    \end{tabular}
    \label{tab:hautaketa}
\end{table}

\textbf{TODO: Osatu taula esperimentu honekin.}

Amaierako hiztegi-tamaina: \textbf{1.001 atributu} (1.000 numeriko + 1 klase).
500 baino handiagoa den baldintza betetzen da.

\subsection{Test-blind prozesua}

Test-blind bermatzeko prozesu hau jarraitu da:

\begin{enumerate}
    \item \texttt{stwv.setInputFormat(train)} -- Hiztegi guztia \textbf{train multzoan
    soilik} ikasten da. Dev eta test-ek ez dute beren hiztegirik sortzen.
    \item \texttt{Filter.useFilter(dev, stwv)} eta \texttt{Filter.useFilter(test, stwv)}
    -- Train-eko hiztegi bera aplikatzen zaie.
    \item \texttt{selector.setInputFormat(trainVec)} -- Atributu-hautaketa ere
    train multzoan soilik ikasten da.
    \item \texttt{Filter.useFilter(devVec, selector)} eta
    \texttt{Filter.useFilter(testVec, selector)} -- Ranking bera aplikatzen zaie.
\end{enumerate}

Prozesu honek bermatzen du train, dev eta test multzoek \textbf{atributu berdinak}
dituztela, \textit{data leakage}-rik gabe.

\newpage
% =============================================
% 5. SAILKATZAILEA: MULTILAYER PERCEPTRON
% =============================================
\section{Sailkatzailea: Multilayer Perceptron}

\subsection{Deskribapen teorikoa}

Multilayer Perceptron (MLP) sare neuronal artifizial bat da, ikasketa gainbegiratua
erabiltzen duena \cite{goodfellow2016, bishop2006}. Sare neuronal sakonagoen oinarria
da eta sailkapen eta erregresio lanetan oso erabilia da.

MLPk hiru geruza mota ditu:

\begin{itemize}
    \item \textbf{Sarrera-geruza (Input Layer):} Bektorizatutako testuaren atributuak
    jasotzen ditu. Kasu honetan, 1.000 atributu numeriko.
    \item \textbf{Geruza ezkutuak (Hidden Layers):} Informazioa prozesatzen duten
    geruzak. Neurona bakoitzak honako kalkulua egiten du:
    \begin{equation}
        a_j = f\!\left(\sum_{i} w_{ij} x_i + b_j\right)
    \end{equation}
    Non $f$ aktibazio-funtzioa den (sigmoid, tanh, ...), $w_{ij}$ pisuak
    eta $b_j$ biasa.
    \item \textbf{Irteera-geruza (Output Layer):} \textit{spam} edo \textit{ham}
    sailkapena ematen du softmax edo sigmoid funtzioa erabiliz.
\end{itemize}

\textbf{Entrenamendua -- Backpropagation:}
Gradient descent metodoa erabiltzen da errorea minimizatzeko:

\begin{equation}
    w_{ij} \leftarrow w_{ij}
    - \eta \frac{\partial E}{\partial w_{ij}}
    + \alpha \Delta w_{ij}^{(t-1)}
\end{equation}

Non $\eta$ learning rate, $E$ errorea (cross-entropy edo MSE) eta
$\alpha$ momentua diren. Momentuak aurreko eguneraketaren zati bat gehitzen
du konbergentzia azkartzeko eta minimum lokaletan ez blokeatzeko.

\textbf{Abantailak:}
\begin{itemize}
    \item Ez-linealtasun konplexuak ikas ditzake
    \item Testu-datu handiekin ondo funtzionatzen du
    \item Backpropagation bidez ikasketa eraginkorra
\end{itemize}

\textbf{Desabantailak:}
\begin{itemize}
    \item Parametro askoren doikuntza behar du (overfitting arriskua)
    \item Konputazio-kostua altua
    \item ``Kutxa beltza'' da -- emaitzak ez dira erraz interpretagarriak
\end{itemize}

\subsection{Parametro sentikorrak}

Weka-ko \texttt{MultilayerPerceptron} klaseak honako parametro sentikorrak ditu:

\begin{table}[H]
    \centering
    \caption{MLP parametro sentikorrak: teoria eta praktika}
    \begin{tabular}{p{2.5cm}p{2.2cm}p{4cm}p{3cm}}
        \toprule
        \textbf{Parametroa} & \textbf{Tarte teorikoa} & \textbf{Interpretazioa} &
        \textbf{Aztertutakoak} \\
        \midrule
        \texttt{learningRate} ($\eta$) & $(0, 1]$ &
        Ikasketaren abiadura. Balio handiak ezegonkortasuna sor dezake;
        txikiak konbergentzia moteltzen du &
        0,001 / 0,003 / 0,005 / 0,010 \\
        \midrule
        \texttt{momentum} ($\alpha$) & $[0, 1)$ &
        Aurreko eguneraketaren eragina. Minimum lokalak ekidin eta
        konbergentzia azkartzen laguntzen du &
        0,2 / 0,4 / 0,6 / 0,8 \\
        \midrule
        \texttt{hiddenLayers} & $\geq 1$ neurona &
        Sarearen konplexutasuna. Gehiegi izateak overfitting ekar dezake;
        gutxiegi, underfitting &
        3 / 5 / 10 / 20 / 30 \\
        \midrule
        \texttt{trainingTime} & $\geq 1$ &
        Entrenamendu-iterazio kopurua (epokak). Gehiegi izateak overfitting &
        100 / 200 / 300 \\
        \bottomrule
    \end{tabular}
    \label{tab:parametroak}
\end{table}

\subsection{Fine-tuning: parametro ekorketa}

\subsubsection{Ebaluazio-eskema eta metrika}

Fine-tuning prozesuan:
\begin{itemize}
    \item \textbf{Ebaluazio-eskema:} Stratified Hold-Out -- train multzoa
    entrenatzeko, dev multzoa ebaluatzeko
    \item \textbf{Datu-sortak:} \texttt{train\_final.arff} (4.137 inst.) eta
    \texttt{dev\_final.arff} (518 inst.)
    \item \textbf{Metrika:} F-Measure (spam) -- precision eta recall batera
    hartzen ditu, klase desorekarentzat egokia
\end{itemize}

\subsubsection{Gorka-ren ekorketa}

Aztertutako balioak: LR $\in \{0{,}001, 0{,}005, 0{,}010\}$,
Mom $\in \{0{,}2, 0{,}3, 0{,}4\}$, Hidden $\in \{3, 5, 10\}$, Epochs = 100.
27 konbinazio guztira.

\begin{table}[H]
    \centering
    \caption{Gorka: fine-tuning emaitzak (hautaketa)}
    \small
    \begin{tabular}{rrlrrrr}
        \toprule
        \textbf{LR} & \textbf{Mom} & \textbf{Hidden} & \textbf{Epochs} &
        \textbf{F-spam} & \textbf{F-ham} & \textbf{WAvg-F} \\
        \midrule
        0,001 & 0,2 & 3  & 100 & 0,9609 & 0,9835 & 0,9770 \\
        0,001 & 0,2 & 10 & 100 & 0,9660 & 0,9855 & 0,9799 \\
        0,005 & 0,2 & 10 & 100 & 0,9803 & 0,9918 & 0,9885 \\
        0,005 & 0,3 & 10 & 100 & 0,9803 & 0,9918 & 0,9885 \\
        \textbf{0,005} & \textbf{0,4} & \textbf{10} & \textbf{100} &
        \textbf{0,9819} & \textbf{0,9925} & \textbf{0,9894} \\
        0,010 & 0,2 & 10 & 100 & 0,9787 & 0,9911 & 0,9875 \\
        0,010 & 0,4 & 10 & 100 & 0,9739 & 0,9890 & 0,9846 \\
        \bottomrule
    \end{tabular}
    \label{tab:finetuning_gorka}
\end{table}

\textbf{Ondorioak:} LR=0,005 onena. Momentum-ak eragin txikia du.
10 neurona 3 eta 5 baino hobea.
\textbf{Eredu onena: LR=0,005 | Mom=0,4 | Hidden=10 | F-spam=0,9819.}

\subsubsection{Andoni-ren ekorketa}

Aztertutako balioak: LR $\in \{0{,}003, 0{,}005\}$,
Mom $\in \{0{,}4, 0{,}6, 0{,}8\}$, Hidden $\in \{10, 20, 30\}$,
Epochs $\in \{100, 200\}$. 36 konbinazio guztira.

\begin{table}[H]
    \centering
    \caption{Andoni: fine-tuning emaitzak (hautaketa)}
    \small
    \begin{tabular}{rrlrrrr}
        \toprule
        \textbf{LR} & \textbf{Mom} & \textbf{Hidden} & \textbf{Epochs} &
        \textbf{F-spam} & \textbf{F-ham} & \textbf{WAvg-F} \\
        \midrule
        0,003 & 0,4 & 10 & 100 & 0,9787 & 0,9911 & 0,9875 \\
        0,003 & 0,6 & 20 & 100 & 0,9819 & 0,9925 & 0,9894 \\
        0,003 & 0,8 & 10 & 100 & 0,9836 & 0,9932 & 0,9904 \\
        0,003 & 0,8 & 20 & 100 & 0,9835 & 0,9932 & 0,9904 \\
        0,003 & 0,8 & 30 & 100 & 0,9835 & 0,9932 & 0,9904 \\
        0,005 & 0,4 & 10 & 100 & 0,9819 & 0,9925 & 0,9894 \\
        \textbf{0,005} & \textbf{0,8} & \textbf{10} & \textbf{100} &
        \textbf{0,9851} & \textbf{0,9939} & \textbf{0,9913} \\
        0,005 & 0,8 & 20 & 100 & 0,9834 & 0,9932 & 0,9904 \\
        \bottomrule
    \end{tabular}
    \label{tab:finetuning_andoni}
\end{table}

\textbf{Ondorioak:} Momentum altuagoak (0,8) emaitza hobeak ematen ditu.
Epoch kopuruak (100 vs 200) eragin txikia du.
\textbf{Eredu onena: LR=0,005 | Mom=0,8 | Hidden=10 | Epochs=100 | F-spam=0,9851.}

\subsubsection{Ibai-ren ekorketa}

Aztertutako balioak: LR $\in \{0{,}003, 0{,}005\}$,
Mom $\in \{0{,}4, 0{,}6, 0{,}8\}$, Hidden $\in \{10, 20, 30\}$,
Epochs $\in \{100, 200\}$. 36 konbinazio guztira.

\begin{table}[H]
    \centering
    \caption{Ibai: fine-tuning emaitzak (hautaketa)}
    \small
    \begin{tabular}{rrlrrrr}
        \toprule
        \textbf{LR} & \textbf{Mom} & \textbf{Hidden} & \textbf{Epochs} &
        \textbf{F-spam} & \textbf{F-ham} & \textbf{WAvg-F} \\
        \midrule
        0,003 & 0,4 & 10 & 100 & 0,8657 & 0,9531 & 0,9278 \\
        0,003 & 0,8 & 10 & 200 & 0,8841 & 0,9579 & 0,9365 \\
        0,005 & 0,4 & 10 & 100 & 0,8263 & 0,9421 & 0,9085 \\
        0,005 & 0,8 & 20 & 100 & 0,8453 & 0,9468 & 0,9174 \\
        \textbf{0,005} & \textbf{0,8} & \textbf{20} & \textbf{200} &
        \textbf{0,8921} & \textbf{0,9604} & \textbf{0,9406} \\
        \bottomrule
    \end{tabular}
    \label{tab:finetuning_ibai}
\end{table}

\textbf{Ondorioak:} Ibai-ren emaitzak nabarmen txikiagoak dira (F-spam $\approx$ 0,89).
Honen arrazoia \texttt{EmailVectorizerAndSelector.java}-ko konfigurazio-aldea
da ziurrenik: \texttt{WordsToKeep} edo \texttt{Ranker}-eko atributu-kopuru ezberdina.
\textbf{Eredu onena: LR=0,005 | Mom=0,8 | Hidden=20 | Epochs=200 | F-spam=0,8921.}

\subsubsection{Emaitzen adierazpen grafikoa}

Irudian LearningRate-ren eragina ikusten da F-spam-ean, Gorka-ren esperimentutik:

\begin{figure}[H]
    \centering
    \begin{tikzpicture}
        \begin{axis}[
            width=10cm, height=6cm,
            xlabel={LearningRate},
            ylabel={F-Measure (spam)},
            xtick={0.001, 0.005, 0.010},
            xticklabels={0,001, 0,005, 0,010},
            ymin=0.95, ymax=1.0,
            ymajorgrids=true,
            legend pos=south east,
            title={LearningRate-ren eragina F-spam-ean (Hidden=10, Epochs=100)}
        ]
        \addplot[color=blue, mark=*, thick] coordinates {
            (0.001, 0.9660)
            (0.005, 0.9819)
            (0.010, 0.9787)
        };
        \addlegendentry{Momentum=0,4}

        \addplot[color=red, mark=square*, thick] coordinates {
            (0.001, 0.9660)
            (0.005, 0.9803)
            (0.010, 0.9755)
        };
        \addlegendentry{Momentum=0,3}

        \addplot[color=green!60!black, mark=triangle*, thick] coordinates {
            (0.001, 0.9660)
            (0.005, 0.9803)
            (0.010, 0.9739)
        };
        \addlegendentry{Momentum=0,2}
        \end{axis}
    \end{tikzpicture}
    \caption{LearningRate-ren eragina F-spam-ean (Gorka-ren fine-tuning, Hidden=10)}
    \label{fig:lr_eragina}
\end{figure}

\begin{figure}[H]
    \centering
    \begin{tikzpicture}
        \begin{axis}[
            width=10cm, height=6cm,
            xlabel={Momentum},
            ylabel={F-Measure (spam)},
            xtick={0.4, 0.6, 0.8},
            xticklabels={0,4, 0,6, 0,8},
            ymin=0.97, ymax=0.99,
            ymajorgrids=true,
            legend pos=south east,
            title={Momentum-aren eragina F-spam-ean (LR=0,005, Epochs=100)}
        ]
        \addplot[color=blue, mark=*, thick] coordinates {
            (0.4, 0.9819)
            (0.6, 0.9771)
            (0.8, 0.9851)
        };
        \addlegendentry{Hidden=10}

        \addplot[color=red, mark=square*, thick] coordinates {
            (0.4, 0.9819)
            (0.6, 0.9803)
            (0.8, 0.9834)
        };
        \addlegendentry{Hidden=20}
        \end{axis}
    \end{tikzpicture}
    \caption{Momentum-aren eragina F-spam-ean (Andoni-ren fine-tuning, LR=0,005)}
    \label{fig:mom_eragina}
\end{figure}

\subsection{Sailkatzaile optimoaren inferentzia}

Hiru fine-tuningen emaitzak konparatuz, \textbf{Andoni-ren eredu onena}
aukeratu da proiektuko eredu optimal gisa, F-spam altuena lortzen duelako:

\begin{table}[H]
    \centering
    \caption{Hiru fine-tuningen eredu onenak -- konparaketa}
    \begin{tabular}{lllllr}
        \toprule
        \textbf{Kidea} & \textbf{LR} & \textbf{Mom} & \textbf{Hidden} &
        \textbf{Epochs} & \textbf{F-spam} \\
        \midrule
        Gorka  & 0,005 & 0,4 & 10 & 100 & 0,9819 \\
        \textbf{Andoni} & \textbf{0,005} & \textbf{0,8} & \textbf{10} &
        \textbf{100} & \textbf{0,9851} \\
        Ibai   & 0,005 & 0,8 & 20 & 200 & 0,8921 \\
        \bottomrule
    \end{tabular}
    \label{tab:eredu_onenak}
\end{table}

\textbf{Eredu optimoa:}
\begin{itemize}
    \item LearningRate: 0,005
    \item Momentum: 0,8
    \item HiddenLayers: 10
    \item Epochs: 100
    \item Entrenatzeko datu-multzoa: \texttt{train\_final.arff} (4.137 instantzia)
    \item Gordetako eredua: \texttt{modelo/mlp.model}
\end{itemize}

\newpage
% =============================================
% 6. KALITATE ESTIMATUA
% =============================================
\section{Aurreikusitako kalitate estimatua}

Eredu finalaren kalitatea bi ebaluazio-eskemaren bidez estimatu da,
\textbf{train eta dev instantziak elkartuz} (4.654 instantzia).
Parametro optimoak: LR=0,005 | Mom=0,2 | Hidden=3 | Epochs=300.

\subsection{5-fold Cross Validation}

5-fold CV-n train+dev multzoa 5 zatitan banatzen da; bakoitzean 4 zatirekin
entrenatzen da eta 1arekin ebaluatzen. Prozesu hau 5 aldiz errepikatzen da.

\begin{table}[H]
    \centering
    \caption{5-fold Cross Validation emaitzak}
    \begin{tabular}{lrrr}
        \toprule
        \textbf{Klasea} & \textbf{Precision} & \textbf{Recall} & \textbf{F-Measure} \\
        \midrule
        Spam  & 0,9419 & 0,9489 & 0,9454 \\
        Ham   & 0,9791 & 0,9761 & 0,9776 \\
        WAvg  & 0,9683 & 0,9682 & 0,9682 \\
        \midrule
        \multicolumn{2}{l}{Accuracy} & \multicolumn{2}{r}{\%96,82} \\
        \bottomrule
    \end{tabular}
    \label{tab:5fcv}
\end{table}

Nahaste-matrizea:
\begin{table}[H]
    \centering
    \caption{5-fold CV nahaste-matrizea}
    \begin{tabular}{rr|l}
        \toprule
        \textbf{a} & \textbf{b} & \\
        \midrule
        1.281 & 69    & a = spam \\
        79    & 3.225 & b = ham  \\
        \bottomrule
    \end{tabular}
\end{table}

69 spam mezu ham bezala sailkatu dira (false negative), eta 79 ham mezu
spam bezala (false positive). Spam-eko false negative-ak kritikoagoak
dira posta-kutxan spam mezu gehiago onartzen direlako.

\subsection{5 Repeated Stratified Hold-Out (70/30)}

10 aldiz errepikatu da 70/30 banaketa estratifikatua, emaitzen egonkortasuna
neurtzeko.

\begin{table}[H]
    \centering
    \caption{5 Repeated Stratified Hold-Out -- iterazio bakoitzeko emaitzak}
    \begin{tabular}{lrrr}
        \toprule
        \textbf{Iterazioa} & \textbf{F-spam} & \textbf{F-ham} & \textbf{WAvg-F} \\
        \midrule
        1 & 0,9455 & 0,9771 & 0,9681 \\
        2 & 0,9573 & 0,9815 & 0,9743 \\
        3 & 0,9606 & 0,9831 & 0,9764 \\
        4 & 0,9368 & 0,9757 & 0,9652 \\
        5 & 0,9399 & 0,9752 & 0,9649 \\
        \midrule
        \textbf{Batezbestekoa} & \textbf{0,9480} & \textbf{0,9785} & \textbf{0,9698} \\
        \bottomrule
    \end{tabular}
    \label{tab:rep_ho_iter}
\end{table}

\begin{table}[H]
    \centering
    \caption{5 Repeated Stratified Hold-Out -- batezbestekoa $\pm$ desbideratze estandarra}
    \begin{tabular}{lrrr}
        \toprule
        \textbf{Klasea} & \textbf{Precision} & \textbf{Recall} & \textbf{F-Measure} \\
        \midrule
        Spam & $0{,}9375 \pm 0{,}0180$ & $0{,}9591 \pm 0{,}0135$ & $0{,}9480 \pm 0{,}0094$ \\
        Ham  & $0{,}9831 \pm 0{,}0056$ & $0{,}9741 \pm 0{,}0069$ & $0{,}9786 \pm 0{,}0032$ \\
        WAvg & --                       & --                       & $0{,}9698 \pm 0{,}0047$ \\
        \bottomrule
    \end{tabular}
    \label{tab:rep_ho}
\end{table}

\subsection{Kalitate estimatuaren interpretazioa}

5-fCV-k eta Repeated HO-k antzeko emaitzak ematen dituzte
(F-spam $\approx$ 0,945--0,948), ereduak ondo orokortzen duela adieraziz.
Desbideratze estandar txikiek ($\pm$0,009) emaitzen egonkortasuna
erakusten dute. Ham klaseko F-Measure altuagoa da datuen desoreka dela eta.

\newpage
% =============================================
% 7. IRAGARPENAK -- ERABILTZAILE ROLA
% =============================================
\section{Iragarpenak: erabiltzailearen paketea}

\subsection{Test-blind ebaluazioa}

\texttt{TestModel.java} fitxategiak eredu finala \texttt{test\_final.arff}
multzoan ebaluatu du (517 instantzia, klase ezezagunak):

\begin{table}[H]
    \centering
    \caption{Test multzoko emaitzak (517 instantzia)}
    \begin{tabular}{lr}
        \toprule
        \textbf{Metrika} & \textbf{Balioa} \\
        \midrule
        Accuracy                    & \%97,87 \\
        F-Measure haztatua (WAvg)   & 0,9787  \\
        \midrule
        Spam Precision              & 0,9633  \\
        Spam Recall                 & 0,9633  \\
        Spam F-Measure              & 0,9633  \\
        \bottomrule
    \end{tabular}
    \label{tab:test}
\end{table}

\begin{table}[H]
    \centering
    \caption{Test multzoko nahaste-matrizea}
    \begin{tabular}{rr|l}
        \toprule
        \textbf{a} & \textbf{b} & \\
        \midrule
        289 & 11 & a = spam \\
        11  & 724 & b = ham \\
        \bottomrule
    \end{tabular}
\end{table}

Test F-spam (0,9633) kalitate estimatutakoa (0,9480) baino altuagoa da,
ereduak benetako datu ezezagunekin ere ondo funtzionatzen duela erakutsiz.

\subsection{Iragarri.java -- Mezu elektroniko berriak sailkatzeko}

\texttt{Iragarri.java} klaseak sarrera direktorio batean dauden \texttt{.txt} fitxategiak (mezu elektronikoak) kargatzen ditu eta bakoitzari dagokion \textit{spam}/\textit{ham} iragarpena egiten dio. Prozesu osoak entrenamenduko pipeline berbera errepikatzen du koherentzia bermatzeko:

\begin{enumerate}
    \item Fitxategiak direktorio batetik kargatu eta \texttt{Instances} objektu bilakatu.
    \item Entrenamenduan gordetako filtroak kargatu eta aplikatu: bektorizazioa (\texttt{vectorizer.ser}) eta atributu-hautaketa (\texttt{selector.ser}).
    \item Entrenatutako eredua kargatu (\texttt{mlp.model}).
    \item \texttt{classifyInstance()} eta \texttt{distributionForInstance()} bidez iragarpenak eta konfidantza-mailak kalkulatu.
    \item Emaitzak fitxategi batean gorde (\texttt{iragarpenak.txt}) eta esperimentu log-ean erregistratu.
\end{enumerate}

Exekuzio adibidea:
\begin{verbatim}
java -cp "lib\weka.jar;bin" Iragarri \
    modelo/mlp.model data_proba/  emaitzak/
\end{verbatim}

Kontsolan eta \texttt{iragarpenak.txt} fitxategian lortzen den irteera:
\begin{verbatim}
==================================================
               IRAGARPENAK - EMAITZAK             
==================================================

Mezua 1: SPAM (konfidantza: 0.9942)
Mezua 2: HAM (konfidantza: 0.8751)

==================================================
Guztira sailkatutako mezuak: 2
==================================================
\end{verbatim}

\newpage
% =============================================
% 8. SOFTWARE DISEINUA
% =============================================
\section{Diseinua}

Helburuetan bi rol bereizten dira: datuen pre-prozesamendua eta klasifikazio/ebaluazioa. Atal bakoitzerako diseinatu dugun softwarearen diseinua grafikoan adierazten da datuen fluxu-diagramen bitartez, non prozesu nagusiak eta datuen transformazioak ikus daitezkeen.

\subsection{Datu-fluxu diagrama}

\begin{figure}[h!]
    \centering
     Hemen jarri dezakezu zuk sortutako irudia, adibidez:
     \includegraphics[width=1.17\textwidth]{diagramas/fuxu_diagrama.png}
    \caption{Datu-fluxu sinplea: sarrera, prozesuak eta irteerak.}
\end{figure}

\subsection{Klase-diagrama}
Gure softwareak 10 Java fitxategik osatzen dute: 3 datuen pre-prozesamendurako (irudian berdez agertzen direnak), 4 ereduaren entrenamendu eta ebaluazioentzat (irudian laranjaz edo urdinez agertzen direnak), eta beste 3 klase lagungarri gisa (irudian morez), exekuzio-fluxu nagusiaren parte ez diren arren, beste klaseek erabiltzen dituztenak.

Hainbat klasek funtzio publikoak dituzte, beste klaseetatik deitzen direnak. Hortaz, fitxategien artean dependentzia erlazio bat sortzen da, \texttt{EmailLoader.java} eta \texttt{ExperimentLogger.java} izanik klase dependente gehien dituzten klaseak.
\begin{figure}[h!]
    \centering
    \includesvg[width=0.95\linewidth]{diagramas/KlaseDiagrama.svg}
    \caption{Klase-diagrama}
    \label{fig:klase-diagrama}
\end{figure}

\subsection{Sarrera eta Irteerak}

\begin{itemize}
    \item \textbf{Sarrera:} ARFF formatuko fitxategiak, adibidez, \texttt{test\_final.arff} edo \texttt{train\_final.arff}, partiketako azpi-multzoak erabiliz.
    \item \textbf{Barneko prozesuak:}
    \begin{itemize}
        \item \textbf{Pre-prozesamendua:} testu normalizazioa (\texttt{TextNormalizer}), tokenizazioa eta atributuen selekzioa (\texttt{StringToWordVector}, \texttt{AttributeSelection}).
        \item \textbf{Inferentzia:} entrenatutako ereduak aplikatzea (\texttt{Classifier} erabiltzen \texttt{Iragarri} eta \texttt{TestModel} klaseetan).
        \item \textbf{Ebaluazioa:} zehaztasuna, recall eta F-measure kalkulatzea (\texttt{Evaluation} klasearen bidez).
    \end{itemize}
    \item \textbf{Irteera:} testuen emaitzak fitxategi testuetan gordetzea (\texttt{.txt}), baita konfusio-matrizeak eta metrika orokorrak.
\end{itemize}

\subsection{Datuen transformazioa}

Datuak sarreratik irteerara prozesu hauen bidez pasatzen dira:

\begin{enumerate}
    \item Sarrera ARFF fitxategia kargatu.
    \item Testu atributuak normalizatu (\texttt{TextNormalizer}).
    \item Tokenizatu eta aukeratu atributuak garrantzitsuenak (\texttt{StringToWordVector}, \texttt{AttributeSelection}).
    \item Aplikatu inferentzia entrenatutako ereduarekin (\texttt{Classifier}).
    \item Ebaluatu emaitzak (\texttt{Evaluation}) eta sortu txostenak (\texttt{.txt} fitxategiak).
\end{enumerate}

Horrela, datuen fluxua eta prozesu nagusiak modu grafiko eta testualean dokumentatzen dira, sarrera eta irteera formatuak argi adieraziz.

% =============================================
% BIBLIOGRAFIA
% =============================================
\newpage
\bibliographystyle{plain}
\begin{thebibliography}{9}

\bibitem{goodfellow2016}
Goodfellow, I., Bengio, Y., \& Courville, A. (2016).
\textit{Deep Learning}. MIT Press.

\bibitem{bishop2006}
Bishop, C. M. (2006).
\textit{Pattern Recognition and Machine Learning}. Springer.

\bibitem{weka}
Hall, M., Frank, E., Holmes, G., Pfahringer, B., Reutemann, P., \& Witten, I. H. (2009).
The WEKA data mining software: An update.
\textit{ACM SIGKDD Explorations Newsletter}, 11(1), 10--18.

\bibitem{manning}
Manning, C. D., Raghavan, P., \& Sch\"utze, H. (2008).
\textit{Introduction to Information Retrieval}. Cambridge University Press.

\bibitem{hastie}
Hastie, T., Tibshirani, R., \& Friedman, J. (2009).
\textit{The Elements of Statistical Learning}. Springer.

\end{thebibliography}

\end{document}